\documentclass[11pt,a4paper]{article}
\usepackage[T1]{fontenc}
\usepackage[utf8]{inputenc}
\usepackage{amsmath,amssymb,amsthm}
\usepackage[margin=1in]{geometry}
\usepackage[colorlinks=true,linkcolor=blue,urlcolor=blue]{hyperref}
\theoremstyle{plain}
\newtheorem{theorem}{Theorem}
\newtheorem{lemma}[theorem]{Lemma}
\newtheorem{proposition}[theorem]{Proposition}
\newtheorem{corollary}[theorem]{Corollary}
\theoremstyle{definition}
\newtheorem{remark}[theorem]{Remark}

\title{The empirical recurrence for OEIS A186798, proved by transfer matrix}
\author{Adrian Perez Fontelles\\ \small Independent researcher}
\date{2 September 2026}
\begin{document}
\maketitle

\begin{abstract}
OEIS A186798 counts the $(n+2)\times3$ arrays over $\{0,\dots,1\}$ in which
neighbouring $3\times3$ subblocks do not commute as matrices, and carries an empirical recurrence
of order $55$ contributed by R.~H.~Hardin. It is true, and it is decidable rather than
empirical. The condition ties together $3+1$ consecutive rows, so $3$ consecutive rows are
a state, the count is a walk count on $S=512$ of them, and the conjectured recurrence is then
settled by a finite exact computation.
\end{abstract}

\noindent\small 2020 Mathematics Subject Classification. 05A15, 15A27, 15B36.\normalsize

\section{The conjecture}

OEIS A186798 is ``Number of (n+2)X3 binary arrays with no 3X3 subblock commuting with any horizontal or vertical neighbor 3X3 subblock.''. It has offset $1$ and begins
\[
512, 3984, 31204, 244558, 1916364, 15016660, 117671951, 922087441,\ \dots
\]
The entry states, as a conjecture contributed by its author and never marked settled:
\begin{quote}
Empirical: a(n)=25*a(n-2)+163*a(n-3)+635*a(n-4)+1829*a(n-5)+4219*a(n-6)+8077*a(n-7)+13098*a(n-8)+18029*a(n-9)+20940*a(n-10)+19779*a(n-11)+13470*a(n-12)+2601*a(n-13)-9874*a(n-14)-18715*a(n-15)-18124*a(n-16)-4854*a(n-17)+16831*a(n-18)+36180*a(n-19)+40250*a(n-20)+25914*a(n-21)+113*a(n-22)-24900*a(n-23)-44479*a(n-24)-58414*a(n-25)-63227*a(n-26)-47878*a(n-27)-6746*a(n-28)+46058*a(n-29)+82994*a(n-30)+83852*a(n-31)+54756*a(n-32)+17658*a(n-33)-7116*a(n-34)-15427*a(n-35)-12634*a(n-36)-6536*a(n-37)-1285*a(n-38)+2283*a(n-39)+5335*a(n-40)+8616*a(n-41)+10232*a(n-42)+8392*a(n-43)+3550*a(n-44)-893*a(n-45)-2547*a(n-46)-1672*a(n-47)-233*a(n-48)+436*a(n-49)+382*a(n-50)+134*a(n-51)+7*a(n-52)-15*a(n-53)-6*a(n-54)-a(n-55) for n>56
\end{quote}
As of the ``Last modified'' line on the live entry (Oct 03 2025 03:42:18, revision 9) this is still
recorded as empirical, and nothing on the entry records it as proved.

\section{What is being counted}

The contiguous $3\times3$ subblocks of an $(n+2)\times3$ array form a grid of $n$
rows and $1$ column. Two subblocks that are neighbours in that
grid OVERLAP: horizontally they share $3\times2$ entries, vertically
$2\times3$. Nevertheless each is a $3\times3$ integer matrix in its own right, and
the entry requires that every neighbouring pair, horizontal or vertical, do not commute --- that is,
$AB=BA$ as matrices over $\mathbf{Z}$.

At $n=1$ the array has $3$ rows and the grid is a single row of $1$
subblock, with no neighbouring pair at all; the entry's first term is
$512$, which the model reproduces along with every later term.

\section{The count is a walk count}

A subblock row is fixed by $3$ consecutive array rows, and the vertical condition compares
two consecutive subblock rows, hence $3+1$ consecutive array rows. So take as state the
$3$-tuple of consecutive rows whose subblock row has just been completed. A step appends one
array row: that fixes the next subblock row, and the horizontal pairs inside it and the
vertical pairs against the previous one are all settled there. An array of $n+2$ rows is
a walk of $n-1$ steps and every walk comes from exactly one array, so
\[
a(n)\;=\;\iota^{\!\top}M^{\,n-1}\tau,\qquad\iota=\tau=(1,\dots,1)^{\!\top},
\]
on the $S=512$ admissible $3$-tuples. In particular $a$ is $C$-finite.

This state does not compress. What a step needs of the past is the previous subblock row, and
for $3\ge3$ that row already determines the $3$ array rows it was built from, so nothing
is lost by carrying the rows and nothing would be gained by carrying the matrices instead. The
size of the construction is therefore $2^{3\cdot3}$ before the condition prunes
it, which is what puts the wider members of this family out of reach rather than settling them.

\section{The proof}

\begin{lemma}\label{lem:crit}
Let $q(t)=t^{55}-\sum_i c_it^{55-i}$ be the characteristic polynomial of the
conjectured recurrence. The residual $a(n)-\sum_ic_ia(n-i)$ equals
$\iota^{\!\top}M^{\,j}q(M)\tau$ for the corresponding walk index $j$, so the recurrence
holds from some point on if and only if $\iota^{\!\top}M^{j}q(M)\tau=0$ for all large $j$,
and it is enough to examine $j<S$.
\end{lemma}

\begin{proof}
Factor $M^{j}$ out of $M^{j+55}-\sum_ic_iM^{j+55-i}$. For the bound, $M^{S}$
is an integer combination of $I,M,\dots,M^{S-1}$ by Cayley--Hamilton, so the numbers
$u_j=\iota^{\!\top}M^{j}q(M)\tau$ obey the monic recurrence given by the characteristic
polynomial of $M$; $S$ consecutive zeros force every later one, and the last nonzero $u_j$
pins the threshold exactly.
\end{proof}

\begin{theorem}
\[
a(n)\;=\;25\,a(n-2) + 163\,a(n-3) + 635\,a(n-4) + 1829\,a(n-5) + 4219\,a(n-6) + 8077\,a(n-7) + 13098\,a(n-8) + 18029\,a(n-9) + 20940\,a(n-10) + 19779\,a(n-11) + 13470\,a(n-12) + 2601\,a(n-13) - 9874\,a(n-14) - 18715\,a(n-15) - 18124\,a(n-16) - 4854\,a(n-17) + 16831\,a(n-18) + 36180\,a(n-19) + 40250\,a(n-20) + 25914\,a(n-21) + 113\,a(n-22) - 24900\,a(n-23) - 44479\,a(n-24) - 58414\,a(n-25) - 63227\,a(n-26) - 47878\,a(n-27) - 6746\,a(n-28) + 46058\,a(n-29) + 82994\,a(n-30) + 83852\,a(n-31) + 54756\,a(n-32) + 17658\,a(n-33) - 7116\,a(n-34) - 15427\,a(n-35) - 12634\,a(n-36) - 6536\,a(n-37) - 1285\,a(n-38) + 2283\,a(n-39) + 5335\,a(n-40) + 8616\,a(n-41) + 10232\,a(n-42) + 8392\,a(n-43) + 3550\,a(n-44) - 893\,a(n-45) - 2547\,a(n-46) - 1672\,a(n-47) - 233\,a(n-48) + 436\,a(n-49) + 382\,a(n-50) + 134\,a(n-51) + 7\,a(n-52) - 15\,a(n-53) - 6\,a(n-54) - a(n-55)
\]
for every $n>56$.
\end{theorem}

\begin{proof}
The vector $w=q(M)\tau$ was formed by $55$ matrix--vector products in exact integer
arithmetic and $\iota^{\!\top}M^{j}w$ evaluated for $j=0,1,\dots$, again exactly; those
integers vanish from the index corresponding to $n=56+1$ onwards and the run continued
until the working vector was identically zero, which settles every larger $j$ at once.
\end{proof}

\section{Verification}

Three checks, each independent of the algebra above.

First, the model reproduces all $17$ terms the entry publishes, exactly, in integer
arithmetic. This is what ties the model to the entry: the digraph is built by reading the
entry's English, and a misreading gives different counts.

Second, the conjectured recurrence was evaluated directly on those published terms, with no
matrices involved, and holds wherever the proved range and the data overlap.

Third, the annihilation test of Lemma~\ref{lem:crit} was carried out in exact integer
arithmetic on the whole vector rather than on a sampled prefix.

\begin{thebibliography}{9}
\bibitem{oeis} The OEIS Foundation, \emph{The On-Line Encyclopedia of Integer
Sequences}, \url{https://oeis.org}, 2026. Sequence A186798.
\bibitem{stanley} R.~P.~Stanley, \emph{Enumerative Combinatorics}, Volume 1, 2nd ed.,
Cambridge University Press, 2011. (Transfer-matrix method, Section 4.7.)
\bibitem{suprunenko} D.~A.~Suprunenko and R.~I.~Tyshkevich, \emph{Commutative Matrices},
Academic Press, 1968.
\bibitem{horn} R.~A.~Horn and C.~R.~Johnson, \emph{Matrix Analysis}, 2nd ed., Cambridge
University Press, 2013.
\end{thebibliography}

\end{document}
